\subsection{可乘性判别法}

以下我们恒设赋范代数 A 是完备的。

定理 1. 设 \((\pmb{x}_n)_{n\in \mathbb{N}}\) 是完备赋范代数 A 中的点的一个序列。a) 若 \((\pmb{x}_n)\) 可乘且其积是 A 的一个单位，则对每个 \(\varepsilon >0\) 存在 N 的有限子集 \(\mathbf{J}_0\) ，使得对 N 的每个不交 \(\mathbf{J}_0\) 的有限子集 L ，有 \(||e - p_{\mathrm{L}}||\leqslant \varepsilon\) 。

\begin{enumerate}
\def\labelenumi{\alph{enumi})}
\setcounter{enumi}{1}
\item
  反之，若序列 \((\pmb{x}_n)\) 满足此条件，则它可乘。此外，若每个 \(\pmb{x}_n\) 是单位，则 \(\prod_{n\in \mathbb{N}}\pmb{x}_n\) 是单位。
\end{enumerate}

a) 设 \(\pmb{p}\) 是可乘序列 \((\pmb{x}_n)\) 的积，且设 \(\pmb{p}\) 是 A 中的单位；则（第 IX 章，§ 3，第 7 号，命题 13）存在 \(\alpha > 0\) 与 \(a > 0\) ，使得对一切满足如下条件的 \(y \in A\) ：

\[
\left| \left| \boldsymbol {y} - \boldsymbol {p} \right| \right| \leqslant \alpha ,
\]

y 是单位且 \(\|y^{-1}\|\leqslant a\) 。由假设，对每个 \(\varepsilon\) （满足 \(o<\varepsilon<\alpha\) ），存在 N 的有限子集 \(H_{0}\) ，使得对 N 的每个包含 \(H_{0}\) 的有限子集 H ，有 \(\|p_{H}-p\|\leqslant\varepsilon\) 。设 \(J_{0}=[0,m]\) 是包含 \(H_{0}\) 的 N 的一个区间；对 N 的每个不交 \(J_{0}\) 的有限子集 L ，属于 L 的整数都大于属于 \(H_{0}\) 的整数；因此，若 \(H=H_{0}\cup L\) ，我们有 \(p_{H}=p_{H_{0}}p_{L}\) 。现在，由于 \(\|p_{H_{0}}-p\|\leqslant\varepsilon\leqslant\alpha\) ，\(p_{H_{0}}\) 是单位，且

\[
\left| \left| e - p _ {\mathrm{H} _ {0}} ^ {- 1} p \right| \right| \leqslant \varepsilon \left| \left| p _ {\mathrm{H} _ {0}} ^ {- 1} p \right| \right| \leqslant a \varepsilon ;
\]

由于 \(||p_{\mathrm{H_0}}p_{\mathrm{L}} - p|| \leqslant \varepsilon\) ，我们推出

\[
\left| \left| p _ {\mathrm{L}} - p _ {\mathrm{H} _ {0}} ^ {- 1} p \right| \right| \leqslant \varepsilon \left| \left| p _ {\mathrm{H} _ {0}} ^ {- 1} \right| \right| \leqslant a \varepsilon ,
\]

而最终 \(||e - p_{\mathrm{L}}|| \leqslant 2a\varepsilon\) 。

\begin{enumerate}
\def\labelenumi{\alph{enumi})}
\setcounter{enumi}{1}
\tightlist
\item
  设对每个 \(\varepsilon > 0\) ，存在有限子集 \(J_0\) （ \(\mathbf{N}\) 的），使得对每个有限子集 \(\mathbf{L}\) （ \(\mathbf{N}\) 的，不交 \(J_0\) ），有 \(\|e - p_{\mathbf{L}}\| \leqslant \varepsilon\) 。设 \(H_0 = [0, p]\) 是 \(\mathbf{N}\) 的一个包含 \(J_0\) 的区间；则每个有限子集 \(\mathbf{H}\) （ \(\mathbf{N}\) 的，包含 \(H_0\) ）可写成 \(H_0 \cup L\) 的形式，其中 \(\mathbf{L}\) 中的整数都大于 \(H_0\) 中的整数；因此我们有 \(p_H = p_{H_0}p_L\) ，且由于 \(\mathbf{L}\) 不交 \(J_0\) ，\(||p_H - p_{H_0}|| \leqslant \varepsilon ||p_{H_0}||\) ，从而 \(||p_H|| = (1 + \varepsilon)||p_{H_0}||\) 。若 \(p_{H_0} = 0\) ，则序列 \((x_n)\) 显然可乘且其积为 0 ；除去这一平凡情形，存在区间 \(H_1 = [0, q]\) （它包含 \(H_0\) ），使得对 N 的每个不交 \(H_1\) 的有限子集 L ，有 \(||e - p_L|| \leqslant \varepsilon (||p_{H_0}||)^{-1}\) 。如上可得，对每个有限子集 \(H \supset H_1\) ，
\end{enumerate}

\[
\left| \left| p _ {\mathrm{H}} - p _ {\mathrm{H} _ {4}} \right| \right| \leqslant \left(\left| \left| p _ {\mathrm{H} _ {0}} \right| \right|\right) ^ {- 1} \left| \left| p _ {\mathrm{H} _ {4}} \right| \right| \varepsilon \leqslant \varepsilon (\mathrm{I} + \varepsilon).
\]

因此柯西判别法表明 \(J \to p_J\) 在 \(A\) 中（关于有序集 \(\mathfrak{F}(N)\) ）有极限。

若所有 \(x_{n}\) 都是单位，则所有有限部分积 \(p_{J}\) 也是单位；从而对每个包含 \(H_{0}\) 的有限子集 H ，我们有

\[
\left| \left| e - p _ {\mathrm{H} _ {0}} ^ {- 1} p _ {\mathrm{H}} \right| \right| \leqslant \varepsilon ,
\]

这表明，在 A 的单位所成的乘法群 G 中，映射 \(J \rightarrow p_{J}\) 下 \(\mathfrak{F}(N)\) 的截口滤的像是关于 G 的左一致结构的一个柯西滤基；但由于 G 完备（第 IX 章，§ 3，第 7 号，命题 13），映射 \(J \rightarrow p_{J}\) 的极限属于 G 。

注. 若 \((\pmb{x}_n)\) 可乘而其积不是单位，定理 1 的条件未必满足：例如，若所有 \(x_n\) 都等于同一元素 \(x\) ，其中 \(||x|| < 1\) ，则序列 \((\pmb{x}_n)\) 可乘且其积为 0 ，而对 N 的每个非空有限子集 H ，有 \(\| p_{\mathbf{H}} \| \leqslant \| x \| < 1 \leqslant \| e \|\) 。

推论 1. 若 \((\pmb{x}_n)\) 是积为 A 的单位的可乘序列，则 \(\lim_{n\to \infty}\pmb{x}_n = \pmb{e}\) 。

推论 2. 若 \((\pmb{x}_n)\) 是积为 A 的单位的可乘序列，则它的每个子序列 \((\pmb{x}_{n_k})_{k\in \mathbb{N}}\) （ \((\pmb{x}_n)[(n_k)\) 为严格递增的整数序列{]}）都是可乘的。

这由定理 1 的判别法立得。

定理 2. 设 A 是完备赋范代数。若 \((\pmb{u}_n)\) 是 A 的元素的一个绝对收敛级数，则序列 \((\pmb{e} + \pmb{u}_n)\) 在 A 中可乘；且若所有元素 \(\pmb{e} + \pmb{u}_n\) 都是 A 中的单位，则 \(\prod_{n\in \mathbb{N}}(\pmb {e} + \pmb{u}_n)\) 也是单位。

我们应用定理 1 的判别法。对 N 的每个有限子集 L ，有 \(p_{L} - e = \prod_{n \in L} (e + u_n) - e = \sum_{M} \left( \prod_{n \in M} u_n \right)\) ，其中 M 跑遍 L 的所有非空子集所成的集（以诱导的序线性排序）。由于 \(\left\| \prod_{n\in \mathbf{M}}\pmb {u}_n\right\| \leqslant \prod_{n\in \mathbf{M}}||\pmb {u}_n||,\) 我们可以写

\[
\left| \left| \boldsymbol {p} _ {\mathrm{L}} - \boldsymbol {e} \right| \right| \leqslant \sum_ {\mathbf {M}} \left(\prod_ {n \in \mathbf {M}} \left| \left| \boldsymbol {u} _ {n} \right|\right)\right) = \prod_ {n \in \mathbf {L}} (\mathrm{I} + \left| \left| \boldsymbol {u} _ {n} \right|\right) - \mathrm{I}.
\]

现在由于通项为 \(||\boldsymbol{u}_n||\) 的级数依假设收敛，序列 \((\mathbf{I} + ||\boldsymbol{u}_n||)\) 在 \(\mathbf{R}_+^*\) 中可乘（第 IV 章，§ 7，第 4 号，定理 4）。因此对每个 \(\varepsilon > 0\) 存在有限子集 \(J_0\) （ \(\mathbf{N}\) 的），使得对每个有限子集 \(\mathbf{L}\) （ \(\mathbf{N}\) 的，不交 \(J_0\) ），有 \(\left| \prod_{n \in \mathbf{L}} (\mathbf{I} + ||\boldsymbol{u}_n||) - \mathbf{I} \right| \leqslant \varepsilon\) ；由此即得结论。

推论. 若通项为 \(u_{n}\) 的级数绝对收敛，且诸元素 \(e + u_{n}\) 都不是 A 中的零因子，则积 \(\prod_{n \in \mathbb{N}} (e + u_{n})\) 不是 A 中的零因子。

只有有限个整数 \(n\) 满足 \(||\boldsymbol{u}_n|| \geqslant 1\) 。设 \(J = [0, m]\) 是 \(N\) 的一个包含所有这些整数的区间。序列 \((\boldsymbol{e} + \boldsymbol{u}_n)\) 的积是 \(p_J\) 与元素 \(\prod_{n > m} (\boldsymbol{e} + \boldsymbol{u}_n)\) 的积，后者的所有因子都是单位（第 IX 章，§ 3，第 7 号，命题 12 的推论），因而它本身是单位；由于 \(p_J\) 是有限个非零因子的积，它不是零因子，因此 \(\prod_{n \in N} (\boldsymbol{e} + \boldsymbol{u}_n)\) 不是零因子。

定理 2 给出的可乘性充分条件一般不是必要的（参看习题 6）。然而，在 A 是域 R 上有限秩代数（即 A 作为 R 上的向量空间是有限维的）这一重要情形中它是必要的；特别地，若 A 是四元数体 H 或矩阵代数 \(\mathbf{M}_n(\mathbf{R})\) ，就是这种情形：

命题 1. 设 A 是 R 上的有限秩赋范代数。若 \((\boldsymbol{e} + \boldsymbol{u}_{n})\) 是 A 中的可乘序列，其积是 A 的单位，则通项为 \(u_{n}\) 的级数绝对收敛。

由第 VII 章，§ 3，第 1 号，命题 2 ，存在数 \(c > 0\) ，使得对每个有限族 \((\pmb{x}_i)_{i \in \mathbf{I}}\) （ \(\mathbf{A}\) 的点组成的），我们有

\[
\sum_ {i \in \mathbf {I}} | | \boldsymbol {x} _ {i} | | \leqslant c. \sup _ {\mathrm{J} \subset \mathbf {I}} \left\| \sum_ {i \in \mathrm{J}} \boldsymbol {x} _ {i} \right\|.\tag{1}
\]

设 \((\pmb{a}_n)_{n\in \mathbb{N}}\) 是 A 的元素的任意序列。对 N 的每个有限子集 I ，令

\[
\boldsymbol {p} _ {\mathbf {I}} = \prod_ {i \in \mathbf {I}} (\boldsymbol {e} + \boldsymbol {a} _ {i}), \quad s _ {\mathbf {I}} = \sum_ {i \in \mathbf {I}} \boldsymbol {a} _ {i}, \quad \sigma_ {\mathbf {I}} = \sum_ {i \in \mathbf {I}} | | \boldsymbol {a} _ {i} | |.
\]

引理 1. 对 N 的每个有限子集 I ，令 \(\varphi(\mathrm{I}) = \sup_{\mathrm{J} \subset \mathrm{I}} ||\boldsymbol{p}_{\mathrm{J}} - \boldsymbol{e}||\) 。则对 I 的每个子集 J 我们有

\[
\left| \left| p _ {J} - e - s _ {J} \right| \right| \leqslant \varphi (I) \sigma_ {J}.\tag{2}
\]

若 \(J\) 为空则引理显然；我们对 \(J\) 的元素个数用归纳法证明。设 \(J = K \cup \{j\}\) ，其中 \(\{j\}\) 严格大于每个 \(i \in K\) 。则 \(p_J = p_K(e + a_j)\) 且

\[
s _ {\mathrm{J}} = s _ {\mathrm{K}} + a _ {j},
\]

从而

\[
p _ {J} - e - s _ {J} = p _ {K} - e - s _ {K} + (p _ {K} - e) a _ {j},
\]

而由归纳假设与 \(\varphi(\mathbf{I})\) 的定义，我们有

\[
\left| \left| p _ {J} - e - s _ {J} \right| \right| \leqslant \varphi (I) \sigma_ {K} + \varphi (I) \left| \left| a _ {j} \right| \right| = \varphi (I) \sigma_ {J},
\]

这就证明了引理。

引理 2. 若 I 是 N 的有限子集，满足 \(\varphi(\mathbf{I}) < \mathbf{i}/c\) ，则

\[
\sigma_ {\mathbf {I}} \leqslant c _ {\varphi} (\mathbf {I}) / (\mathbf {I} - c _ {\varphi} (\mathbf {I})).
\]

因为 \(\sigma_{\mathbf{J}} \leqslant \sigma_{\mathbf{I}}\) 对每个子集 \(J\) （ \(I\) 的）成立，由 (2) 得

\[
\left| \left| s _ {J} \right| \right| \leqslant \varphi (I) \sigma_ {I} + \left| \left| p _ {J} - e \right| \right| \leqslant (I + \sigma_ {I}) \varphi (I);
\]

又由于 \(\sigma_{\mathbf{I}} \leqslant c.\sup_{\mathbf{J} \in \mathbf{I}} ||s_{\mathbf{J}}||\) （由 (1) ），可得

\[
\sigma_ {\mathbf {I}} \leqslant c \varphi (\mathbf {I}) (\mathbf {I} + \sigma_ {\mathbf {I}}),
\]

由此即得结果。

现在设 \((e + u_{n})\) 是 A 中的可乘序列，其积为单位；由定理 1 存在 N 的有限子集 \(J_{0}\) ，使得对 N 的每个不交 \(J_{0}\) 的有限子集 H ，我们有

\[
\left\| \prod_ {i \in \mathbf {H}} (e + u _ {i}) - e \right\| \leqslant 1 / 2 c.
\]

由引理 2 得 \(\sum_{i\in\mathbf{H}}||\boldsymbol{u}_{i}||-\mathrm{i}\) （对 N 的每个不交 \(J_{0}\) 的有限子集 H ），从而（第 IV 章，§ 7，第 1 号，定理 1）族 \((||\boldsymbol{u}_{n}||)\) 在 R 中可和。

